\section{A family of counterexamples in type D}\label{sec:type-d}

We prove Theorem~\ref{thm:mainD} by an incidence construction independent of torsion detection and all finite certificates. Fix $r\ge k\ge3$ and put $p=r-k+2$. Let $V$ be a split quadratic space of dimension $2r+4$ with isotropic flag $E_\bullet$. Set $R=E_r$, $F_0=R^\perp$, $P=E_p$, and $H_0=F_0/R$. The quadratic space $H_0$ has dimension four. Choose the ruling of maximal isotropic planes containing $E_{r+2}/R$ and denote its parameter space by $B_0\simeq\PP^1$.

Identify $H_0$ with $U_2\otimes W_2$, with determinant quadratic form. The planes of the chosen ruling are $A_x=\C x\otimes W_2$ and their universal bundle is $\mathcal A\simeq\OO_{\PP^1}(-1)\otimes W_2$. Let $\mathcal E$ be its inverse image in $F_0\otimes\OO_{B_0}$. Choosing a splitting gives $\mathcal E\simeq(R\otimes\OO)\oplus\mathcal A$, of rank $r+2$ and determinant $\OO(-2)$. Every fibre $\mathcal E_A$ is isotropic in $V$.

\subsection{The incidence source and its Schubert image}
On $Z=\Gr_{B_0}(k,\mathcal E)$, let $\mathcal L$ be the tautological subbundle and $\mathcal Q=\mathcal E/\mathcal L$, of rank $p$. Define
\begin{equation}\label{eq:D-incidence}
 \begin{gathered}
 I=\{(A,L)\in Z:\det(P\longrightarrow\mathcal Q)=0\},
 \\ q:I\longrightarrow\OG(k,V),\quad (A,L)\longmapsto L,
 \end{gathered}
\end{equation}
and let $X=q(I)_{\red}$.

\begin{lemma}\label{D:source}
The variety $I$ is integral, normal, and Gorenstein, of dimension $kp$. Its singular locus is the irreducible codimension-three locus $I_2=\{(A,L):\dim(L\cap P)\ge2\}$. Its generic transverse singularity is the node $ad-bc=0$.
\end{lemma}
\begin{proof}
Choose a line $K\subset P$ and then a $k$-plane $L\subset\mathcal E_A$ containing it. This is an integral Grassmann bundle over $\PP(P)\times B_0$, mapping onto the support of $I$. The open set with $\dim(L\cap P)=1$ is nonempty and has a unique such line, so this support is irreducible. A corank-one determinant chart has a nonzero linear parameter, proving generic reducedness. The nonzero determinant section cuts out a Cartier hypersurface in the smooth integral variety $Z$. It is Cohen--Macaulay with no embedded associated primes, so generic reducedness and irreducible support make it integral. Since $\dim Z=kp+1$, its dimension is $kp$.

At a point with $\dim(L\cap P)=c$, the map $P\to\mathcal Q$ has corank $c$. Splitting off an invertible block in an ordinary Grassmann chart leaves the determinant of a matrix of $c^2$ independent parameters; the other coordinates are smooth parameters. One can see the independence directly by decomposing the ambient vector space into $L\cap P$, complements in $L$ and $P$, and a complement of $L+P$. The graph maps from $L\cap P$ to the last complement give precisely that $c\times c$ block. The singular locus is therefore exactly $c\ge2$ \cite{BV88,BH93}.

Choosing a two-plane in $P$ and then a $k$-plane containing it gives an integral incidence variety onto $I_2$. The generic intersection is two-dimensional, where the chosen two-plane is unique, and the dimension is $1+2(p-2)+(k-2)p=kp-3$. Thus $I_2$ is irreducible of codimension three. The hypersurface is regular in codimension one, so Serre's criterion gives normality. It is Gorenstein as a hypersurface. At the generic point of $I_2$, the independent block has size $2\times2$, giving the claimed node.
\end{proof}

\begin{lemma}\label{D:image}
The image $X$ is a Schubert variety, and $q:I\to X$ is proper birational. For $T=\Gr(k,R)$,
\begin{equation}\label{eq:D-small}
 \begin{gathered}
 q^{-1}(T)=T\times\PP^1,\qquad\codim_Iq^{-1}(T)=2k-1,\\
 I\setminus q^{-1}(T)\xrightarrow{\ \sim\ }X\setminus T.
 \end{gathered}
\end{equation}
In particular, $q$ has no exceptional prime divisors.
\end{lemma}
\begin{proof}
The map is proper. Its image is irreducible, closed, and $B$-stable: the chosen Borel preserves $P$, $R$, and $F_0$, and its connectedness makes it preserve the chosen ruling. The orthogonal Grassmannian has finitely many $B$-orbits. An irreducible closed $B$-stable subvariety is therefore the closure of a single orbit, so $X$ is a Schubert variety. In particular, it is normal and Cohen--Macaulay \cite{Ram85,Bri05}.

Let $\pi:F_0\to H_0$ be the quotient. If $L\not\subset R$, its image $\pi L$ is a nonzero isotropic line or plane. It has a unique maximal isotropic extension in the specified ruling. This inverse construction is a morphism: write a nonzero residual vector as a rank-one $2\times2$ matrix and recover its left factor from a nonzero column. These formulae are regular on their nonzero-column opens. For any two residual vectors the left factors agree, because their span belongs to one ruling plane; hence the formulae glue. They give an inverse on the reduced schemes, proving the open isomorphism in \eqref{eq:D-small}.

If $L\subset R$, every ruling plane is allowed, and $\dim(L\cap P)\ge k+p-r=2$. The determinant condition then vanishes identically on the product $T\times B_0$, so the fibre identity is scheme-theoretic. Since $\dim T=k(r-k)$, the exceptional locus has codimension $kp-k(r-k)-1=2k-1\ge5$ in $I$. The open isomorphism proves birationality and absence of exceptional prime divisors.
\end{proof}

For an explicit Weyl-group label, set $n=r+2$ and use positive roots $e_i-e_j,e_i+e_j$ for $i<j$, with simple roots $\alpha_i=e_i-e_{i+1}$ for $i<n$ and $\alpha_n=e_{n-1}+e_n$. Choose a compatible hyperbolic basis so that the chosen ruling contains the residual plane labelled by $-n,-(n-1)$. In signed one-line notation, the minimal parabolic representative is
\begin{equation}\label{eq:D-word}
 w=(p,p+2,p+3,\ldots,r,-n,-(n-1),1,2,\ldots,p-1,p+1),
\end{equation}
omitting empty strings. The first positive block contributes $(k-2)p-1$ inversions with the final positive block; the two negative entries contribute $2p$ negative images of difference roots and one negative image of a sum root. All other contributions vanish. Thus $\ell(w)=kp$, and its unique right descent is $k$. The corresponding fixed $k$-plane belongs to the incidence image, so its Schubert closure, of dimension $kp$, is $X$. The geometric proof of Lemma~\ref{D:image} does not require this label.

\subsection{The canonical obstruction along the contracted locus}
Use additive line-bundle notation and set $H=\det\mathcal L^*$ and $\zeta=\OO_{B_0}(1)$. The relative Grassmann tangent bundle is $\mathcal L^*\otimes\mathcal Q$. Since $\det\mathcal Q=H-2\zeta$, taking its determinant and then adding $K_{B_0}=-2\zeta$ gives $K_Z=-(r+2)H+(2k-2)\zeta$. The determinant section in \eqref{eq:D-incidence} has line bundle $\det\mathcal Q\otimes(\det P)^*=H-2\zeta$. Adjunction therefore yields
\[
 K_I=-(r+1)H+(2k-4)\zeta.
\]
Here adjunction applies to a Cartier hypersurface in a smooth variety, even at the singular points of the hypersurface \cite[Chapter~II, Section~8]{Har77}\cite[Chapter~3]{BH93}.

\begin{proposition}\label{D:nonQ}
The ambient variety $X$ is non-$\Q$-Gorenstein at every scheme point of $T$ and is Gorenstein on $X\setminus T$.
\end{proposition}
\begin{proof}
The complement is isomorphic to an open subset of the Gorenstein hypersurface $I$. For $x\in T$, its fibre curve is $C_x\simeq\PP^1_{\kappa(x)}$, and $H$ and $\zeta$ have degrees zero and one on it. Thus
\[
 \deg(K_I|_{C_x})=2k-4>0.
\]

Suppose $mK_X$ were Cartier near $x$ for some $m>0$. Choose canonical Weil divisors using the same rational top form on the common function field of $I$ and $X$. The map is an isomorphism outside subsets of codimension at least two in both varieties and has no exceptional prime divisor. Hence $mK_I=q^*(mK_X)$ as Weil divisors over this neighbourhood. Both sides are Cartier on the normal variety $I$, so their line bundles agree. The right side has degree zero on $C_x$, while the left side has degree $m(2k-4)>0$, a contradiction. This argument works over $\kappa(x)$ and therefore covers nonclosed points as well.
\end{proof}

\begin{proof}[Proof of Theorem~\ref{thm:mainD}]
Set $S=q(I_2)$. It is irreducible. The exceptional locus has codimension at least five, whereas $I_2$ has codimension three. Its generic point is therefore in the isomorphism locus, so $S$ has codimension three and generic transverse node. The inequality $\dim(L\cap P)\ge2$ for $L\in T$ gives $T\subset S$. Outside $S$, the map identifies $X$ with a smooth open subset of $I$. The dense subset $S\setminus T$ is singular in $X$, and closedness of the singular locus then gives $\Sing X=S$.

Proposition~\ref{D:nonQ} gives $\NGor(X)=\NonQGor(X)=T$, with $\codim_XT=2k$. The generic transverse node along $S$ is Gorenstein, so $\Bbad(X)=\varnothing$.

The inverse image of $X$ under $G/B\to G/P$ is a complete-flag Schubert variety. This projection is smooth and locally trivial with geometrically irreducible flag variety fibre. The singular, Gorenstein, and $\Q$-Gorenstein loci pull back, and codimensions are unchanged. For the last assertion, the canonical sheaf on a local product is the pullback of the base canonical sheaf tensored with an invertible relative canonical sheaf. Reflexive powers commute with the flat extension, and invertibility is reflected by faithful flatness. Complete-flag counterexamples follow.
\end{proof}

\subsection{An independent toric model at a general point of the locus}
At $L_0\in T$ in the open position $\dim(L_0\cap P)=2$, choose decompositions $R=L_0\oplus W$ and $P=(L_0\cap P)\oplus W$. A Grassmann graph chart writes $L=\operatorname{graph}(\alpha,M)$, where $\alpha:L_0\to W$ and $M:L_0\to H_0$. The $k(r-k)$ coordinates of $\alpha$ are unconstrained. The ruling condition writes $M=x\otimes B$, with $x\in\C^2$ and $B$ a $2\times k$ matrix. Choosing the first two basis vectors in $L_0\cap P$, the incidence condition is $\det B_{\{1,2\}}=0$. Thus the actual chart is
\begin{equation}\label{eq:D-toric-chart}
 \A^{k(r-k)}\times C_k,\qquad
 C_k=\Spec\!\left(\left(\C[x_1,x_2,B_{bi}]/(\det B_{\{1,2\}})\right)^{\Gm}\right),
\end{equation}
where $x$ has weight $+1$ and $B$ has weight $-1$.

To check the scheme identity, the ring before invariants is a rank-four quadric ring times a polynomial ring, hence is normal and integral. Its invariant ring is likewise normal and integral. Pairing positive and negative factors shows that it is generated by the entries $x_aB_{bi}$, so it defines a closed reduced subscheme in the graph-coordinate affine space. Its closed-point support is the incidence image: a nonzero residual map recovers the unique ruling line, and the zero map is represented as well. Both this subscheme and the chart of $X$ are reduced, so equality of their closed-point supports over $\C$ gives equality of defining ideals. The scheme-theoretic identification does not require the equations $x_ax_c\det B_{\{1,2\}}$ to generate the entire defining ideal.

Equivalently, $C_k$ is the full $\OO(1,1)$ section cone over $\PP^1\times Q_{2k-2}$, where $Q_{2k-2}\subset\PP^{2k-1}$ is the rank-four quadric $\det B_{\{1,2\}}=0$. Its degree-$d$ dimension is
\[
 (d+1)\left[\binom{d+2k-1}{2k-1}-\binom{d+2k-3}{2k-1}\right],
\]
with binomial coefficients zero when appropriate. Applying $1+t\frac{d}{dt}$ to $(1+t)/(1-t)^{2k-1}$ gives
\[
 H_{C_k}(t)=\frac{1+2kt+(2k-3)t^2}{(1-t)^{2k}}.
\]
This is an identity for every $k$, not an extrapolation from examples.

Factor the first two columns of $B$ as $v\otimes w$, and call the other $2k-4$ entries $d_j$. The toric monoid has nonnegative coordinates $(u_1,u_2,v_1,v_2,w_1,w_2,d_1,\ldots,d_{2k-4})$, subject to
\[
 V=W,\qquad U=V+D,
\]
where capital letters denote coordinate sums. It is generated by the eight monomials $u_av_bw_c$ and the $4k-8$ monomials $u_ad_j$. Pairing factors in these balance equations proves that this is exactly the nonnegative cone intersected with the indicated lattice, so the monoid is saturated. This lattice is also the group generated by the monoid: the solution $u_1=u_2=k-1$ with all $v_i,w_i,d_j$ equal to one is strictly positive, and adding a sufficiently large multiple of it expresses any lattice vector as a difference of two nonnegative solutions. Each coordinate inequality is a primitive facet: its zero face contains a point with all other coordinates positive, and one of the generators has coordinate value one. The toric canonical divisor is $K=-\sum_\rho D_\rho$ \cite[Theorem~8.2.3]{CLS11}.

If $mK_{C_k}$ were Cartier at the vertex, its divisorial module would have a one-dimensional quotient by the vertex ideal. That quotient is torus-stable. Lift a weight vector using diagonalizability, then apply the positive-degree graded Nakayama lemma to obtain a generator of the whole module. In the full torus grading each character space has dimension at most one, so this generator is a monomial. Ordinary cone homogeneity alone would not justify this last assertion. At every invariant prime divisor, the valuation threshold of $\OO(mK)$ is $m$, so the primitive-facet value of the generator is exactly $m$. All these divisors pass through the vertex. Substitution into the second balance equation forces $2m=(2k-2)m$, impossible for $k\ge3$ and $m>0$.

This independently recovers the infinite-order canonical obstruction. Under the lattice map with coordinates $(V-W,U-V-D)$, the canonical vector maps to $(0,2k-4)$, in agreement with the adjunction calculation. For $(r,k)=(3,3)$, the local six-dimensional cone has Hilbert numerator $(1,6,3)$. For $(r,k)=(5,4)$, the transverse eight-dimensional cone along $T\simeq\PP^4$ has numerator $(1,8,5)$.
